% !TEX root = main.tex
\documentclass[main.tex]{subfiles}
\begin{document}

\section{Background}
\label{sec:background}

% We start by reviewing key concepts required by our contributions.
Before presenting our contributions, we review key concepts underlying our work.
Here, we consider a bounded and simply connected domain $\domain$ in $\smash{\R^d}$.
We define a \emph{partition} of $\domain$ as a collection of $n$ cells $\smash{\cell \!=\! \{\cell_i\}_{i=1}^n}$ of positive volume that tile $\domain$.
We denote the {volume} and {centroid} for each cell $\cell_i$ as
\begin{subequations}
  \label{eq:cell-quantities}
  \begin{empheq}[left={\empheqlbrace\,}]{align}
    \cellvol_i &= \int_{\cell_i} d\xx , \label{eq:cell-volume} \\
    \centroid_i &= \frac{1}{\cellvol_i} \int_{\cell_i} \xx \, d\xx . \label{eq:cell-centroid}
  \end{empheq}
\end{subequations}
% 
Because the cells $\cell$ tile $\domain$, their volumes and centroids must verify two conservation identities:
\begin{subequations}
  \label{eq:partition-feasibility}
  \begin{empheq}[left={\empheqlbrace\,}]{align}
    \sum\nolimits_{i} \cellvol_i 
    &= \int_{\domain} \, d\xx 
    = \domainVol, 
    \label{eq:partition-feasibility-vol} 
    \\
    \sum\nolimits_{i} \cellvol_i \, \centroid_i 
    &= \int_{\domain} \xx \, d\xx 
    = \domainVol \, \centroid_\domain, 
    \label{eq:partition-feasibility-centroid}
  \end{empheq}
\end{subequations}
% 
where $\domainVol$ and $\centroid_\domain$ are the volume and the centroid of $\domain$, respectively.

\paragraph{Power diagram.}
We associate each cell $\cell_i$ with a \emph{weighted site} encoded by a site position $\smash{\qq_i \!\in\! \R^d}$ and a scalar weight $\smash{\weight_i \!\in\! \R}$, which we collect as $\qq \!=\! \{\qq_i\}_{i=1}^n$ and $\weight \!=\! \{\weight_i\}_{i=1}^n$.
We also assume the sites are pairwise distinct.
% : two coinciding sites with equal weights have identical power distances and leave the split of their common cell undetermined, while a coinciding site with a smaller weight has an empty cell.
Computing the \emph{power distance} from a point $\xx$ to a weighted site $i$ as $\smash{\pdist{\xx}{i}}$, we define the \emph{power diagram} of $(\qq,\weight)$ as the partition of $\domain$ that assigns each point $\smash{\xx \!\in\! \domain}$ to the weighted site of smallest power distance.
The cell $\cell_i$ of the resulting power diagram for each $(\qq_i, \weight_i)$ is thus expressed as
\begin{equation}
  \cell_i = \bigl\{ \xx \in \domain : \pdist{\xx}{i} \le \pdist{\xx}{j} \ \ \forall j \neq i \bigr\} .
  \label{eq:power-cell}
\end{equation}
Since the quadratic terms $\smash{\|\xx\|^2}$ in \eqref{eq:power-cell} cancel, the points of equal power distance to two sites form a hyperplane and, consequently, every (interior) cell $\cell_i$ spans a convex polytope formed by the intersection of the corresponding half-spaces \cite{Aurenhammer1987PowerDiagrams}.
We call two cells $\cell_i$ and $\cell_j$ {neighbors}, symbolized as $i \!\sim\! j$, when they share a common facet, which we denote $\face_{ij}$ of measure $\area_{ij}$. %, and which lies on the hyperplane of equal power distance to the sites $i$ and $j$.
For instance, a cell in 2D is a polygon whose $\cellvol_i$ is an {area} and a facet is a line segment whose measure $\area_{ij}$ is a {length}, while a cell in 3D is a polyhedron whose $\cellvol_i$ is a {volume} and a facet has measure $\area_{ij}$ equal to an {area}.
Note also that the power diagram reduces to the classical Voronoi diagram when all weights are equal.

\paragraph{Homothety invariance.}
Since we assume distinct sites and cells with positive volume, any given set of weighted sites $(\qq,\weight)$ defines a unique power diagram that depends smoothly on $(\qq,\weight)$ \cite[Prop.~1.1]{deGournay2019Differentiation}.
The converse, however, does not hold, 
\begin{wrapfigure}[7]{r}{0.36\columnwidth}
  \centering
  \vspace{-1mm}
  \resizebox{\linewidth}{!}{%% generated by figures/nullspace_insets.py; do not edit by hand
\begin{tikzpicture}[x={(1.1cm,-0.15cm)}, y={(0.45cm,0.5cm)}, z={(0cm,0.75cm)}]
  \fill[nsBlue!18] (2.000,0.825,0) -- (2.000,2.000,0) -- (1.131,2.000,0) -- (0.881,0.917,0) -- (0.888,0.911,0) -- (1.983,0.820,0) -- cycle;
  \fill[nsOrange!18] (0.881,0.917,0) -- (1.131,2.000,0) -- (0.000,2.000,0) -- (0.000,0.615,0) -- cycle;
  \fill[nsGreen!18] (1.200,0.000,0) -- (2.000,0.000,0) -- (2.000,0.825,0) -- (1.983,0.820,0) -- (1.233,0.308,0) -- cycle;
  \fill[nsPurple!18] (1.233,0.308,0) -- (1.983,0.820,0) -- (0.888,0.911,0) -- cycle;
  \fill[nsRed!18] (0.000,0.000,0) -- (1.200,0.000,0) -- (1.233,0.308,0) -- (0.888,0.911,0) -- (0.881,0.917,0) -- (0.000,0.615,0) -- cycle;
  \draw[nsFacet] (1.131,2.000,0) -- (0.881,0.917,0);
  \draw[nsFacet] (2.000,0.825,0) -- (1.983,0.820,0);
  \draw[nsFacet] (0.888,0.911,0) -- (1.983,0.820,0);
  \draw[nsFacet] (0.881,0.917,0) -- (0.888,0.911,0);
  \draw[nsFacet] (0.881,0.917,0) -- (0.000,0.615,0);
  \draw[nsFacet] (1.983,0.820,0) -- (1.233,0.308,0);
  \draw[nsFacet] (1.200,0.000,0) -- (1.233,0.308,0);
  \draw[nsFacet] (1.233,0.308,0) -- (0.888,0.911,0);
  \draw[nsDomain] (0.000,0.000,0) -- (2.000,0.000,0) -- (2.000,2.000,0) -- (0.000,2.000,0) -- cycle;
  \filldraw[nsBlue!45, draw=nsBlue, line width=0.4pt] (1.614,1.160,0) -- (1.612,1.189,0) -- (1.606,1.218,0) -- (1.597,1.246,0) -- (1.584,1.272,0) -- (1.567,1.296,0) -- (1.548,1.318,0) -- (1.526,1.337,0) -- (1.502,1.354,0) -- (1.476,1.367,0) -- (1.448,1.376,0) -- (1.419,1.382,0) -- (1.390,1.384,0) -- (1.361,1.382,0) -- (1.332,1.376,0) -- (1.304,1.367,0) -- (1.278,1.354,0) -- (1.254,1.337,0) -- (1.232,1.318,0) -- (1.213,1.296,0) -- (1.196,1.272,0) -- (1.183,1.246,0) -- (1.174,1.218,0) -- (1.168,1.189,0) -- (1.166,1.160,0) -- (1.168,1.131,0) -- (1.174,1.102,0) -- (1.183,1.074,0) -- (1.196,1.048,0) -- (1.213,1.024,0) -- (1.232,1.002,0) -- (1.254,0.983,0) -- (1.278,0.966,0) -- (1.304,0.953,0) -- (1.332,0.944,0) -- (1.361,0.938,0) -- (1.390,0.936,0) -- (1.419,0.938,0) -- (1.448,0.944,0) -- (1.476,0.953,0) -- (1.502,0.966,0) -- (1.526,0.983,0) -- (1.548,1.002,0) -- (1.567,1.024,0) -- (1.584,1.048,0) -- (1.597,1.074,0) -- (1.606,1.102,0) -- (1.612,1.131,0) -- cycle;
  \node[circle, fill=black, inner sep=0.7pt] at (1.390,1.160,0) {};
  \filldraw[nsOrange!45, draw=nsOrange, line width=0.4pt] (0.814,1.380,0) -- (0.811,1.429,0) -- (0.801,1.477,0) -- (0.786,1.523,0) -- (0.764,1.567,0) -- (0.737,1.608,0) -- (0.705,1.645,0) -- (0.668,1.677,0) -- (0.627,1.704,0) -- (0.583,1.726,0) -- (0.537,1.741,0) -- (0.489,1.751,0) -- (0.440,1.754,0) -- (0.391,1.751,0) -- (0.343,1.741,0) -- (0.297,1.726,0) -- (0.253,1.704,0) -- (0.212,1.677,0) -- (0.175,1.645,0) -- (0.143,1.608,0) -- (0.116,1.567,0) -- (0.094,1.523,0) -- (0.079,1.477,0) -- (0.069,1.429,0) -- (0.066,1.380,0) -- (0.069,1.331,0) -- (0.079,1.283,0) -- (0.094,1.237,0) -- (0.116,1.193,0) -- (0.143,1.152,0) -- (0.175,1.115,0) -- (0.212,1.083,0) -- (0.253,1.056,0) -- (0.297,1.034,0) -- (0.343,1.019,0) -- (0.391,1.009,0) -- (0.440,1.006,0) -- (0.489,1.009,0) -- (0.537,1.019,0) -- (0.583,1.034,0) -- (0.627,1.056,0) -- (0.668,1.083,0) -- (0.705,1.115,0) -- (0.737,1.152,0) -- (0.764,1.193,0) -- (0.786,1.237,0) -- (0.801,1.283,0) -- (0.811,1.331,0) -- cycle;
  \node[circle, fill=black, inner sep=0.7pt] at (0.440,1.380,0) {};
  \filldraw[nsGreen!45, draw=nsGreen, line width=0.4pt] (1.730,0.270,0) -- (1.729,0.283,0) -- (1.727,0.296,0) -- (1.722,0.308,0) -- (1.717,0.320,0) -- (1.709,0.331,0) -- (1.701,0.341,0) -- (1.691,0.349,0) -- (1.680,0.357,0) -- (1.668,0.362,0) -- (1.656,0.367,0) -- (1.643,0.369,0) -- (1.630,0.370,0) -- (1.617,0.369,0) -- (1.604,0.367,0) -- (1.592,0.362,0) -- (1.580,0.357,0) -- (1.569,0.349,0) -- (1.559,0.341,0) -- (1.551,0.331,0) -- (1.543,0.320,0) -- (1.538,0.308,0) -- (1.533,0.296,0) -- (1.531,0.283,0) -- (1.530,0.270,0) -- (1.531,0.257,0) -- (1.533,0.244,0) -- (1.538,0.232,0) -- (1.543,0.220,0) -- (1.551,0.209,0) -- (1.559,0.199,0) -- (1.569,0.191,0) -- (1.580,0.183,0) -- (1.592,0.178,0) -- (1.604,0.173,0) -- (1.617,0.171,0) -- (1.630,0.170,0) -- (1.643,0.171,0) -- (1.656,0.173,0) -- (1.668,0.178,0) -- (1.680,0.183,0) -- (1.691,0.191,0) -- (1.701,0.199,0) -- (1.709,0.209,0) -- (1.717,0.220,0) -- (1.722,0.232,0) -- (1.727,0.244,0) -- (1.729,0.257,0) -- cycle;
  \node[circle, fill=black, inner sep=0.7pt] at (1.630,0.270,0) {};
  \filldraw[nsPurple!45, draw=nsPurple, line width=0.4pt] (1.405,0.680,0) -- (1.404,0.687,0) -- (1.403,0.694,0) -- (1.401,0.701,0) -- (1.397,0.707,0) -- (1.393,0.713,0) -- (1.389,0.719,0) -- (1.383,0.723,0) -- (1.377,0.727,0) -- (1.371,0.731,0) -- (1.364,0.733,0) -- (1.357,0.734,0) -- (1.350,0.735,0) -- (1.343,0.734,0) -- (1.336,0.733,0) -- (1.329,0.731,0) -- (1.323,0.727,0) -- (1.317,0.723,0) -- (1.311,0.719,0) -- (1.307,0.713,0) -- (1.303,0.707,0) -- (1.299,0.701,0) -- (1.297,0.694,0) -- (1.296,0.687,0) -- (1.295,0.680,0) -- (1.296,0.673,0) -- (1.297,0.666,0) -- (1.299,0.659,0) -- (1.303,0.653,0) -- (1.307,0.647,0) -- (1.311,0.641,0) -- (1.317,0.637,0) -- (1.323,0.633,0) -- (1.329,0.629,0) -- (1.336,0.627,0) -- (1.343,0.626,0) -- (1.350,0.625,0) -- (1.357,0.626,0) -- (1.364,0.627,0) -- (1.371,0.629,0) -- (1.377,0.633,0) -- (1.383,0.637,0) -- (1.389,0.641,0) -- (1.393,0.647,0) -- (1.397,0.653,0) -- (1.401,0.659,0) -- (1.403,0.666,0) -- (1.404,0.673,0) -- cycle;
  \node[circle, fill=black, inner sep=0.7pt] at (1.350,0.680,0) {};
  \filldraw[nsRed!45, draw=nsRed, line width=0.4pt] (1.014,0.360,0) -- (1.012,0.389,0) -- (1.006,0.418,0) -- (0.997,0.446,0) -- (0.984,0.472,0) -- (0.967,0.496,0) -- (0.948,0.518,0) -- (0.926,0.537,0) -- (0.902,0.554,0) -- (0.876,0.567,0) -- (0.848,0.576,0) -- (0.819,0.582,0) -- (0.790,0.584,0) -- (0.761,0.582,0) -- (0.732,0.576,0) -- (0.704,0.567,0) -- (0.678,0.554,0) -- (0.654,0.537,0) -- (0.632,0.518,0) -- (0.613,0.496,0) -- (0.596,0.472,0) -- (0.583,0.446,0) -- (0.574,0.418,0) -- (0.568,0.389,0) -- (0.566,0.360,0) -- (0.568,0.331,0) -- (0.574,0.302,0) -- (0.583,0.274,0) -- (0.596,0.248,0) -- (0.613,0.224,0) -- (0.632,0.202,0) -- (0.654,0.183,0) -- (0.678,0.166,0) -- (0.704,0.153,0) -- (0.732,0.144,0) -- (0.761,0.138,0) -- (0.790,0.136,0) -- (0.819,0.138,0) -- (0.848,0.144,0) -- (0.876,0.153,0) -- (0.902,0.166,0) -- (0.926,0.183,0) -- (0.948,0.202,0) -- (0.967,0.224,0) -- (0.984,0.248,0) -- (0.997,0.274,0) -- (1.006,0.302,0) -- (1.012,0.331,0) -- cycle;
  \node[circle, fill=black, inner sep=0.7pt] at (0.790,0.360,0) {};
  \draw[nsGuide] (1.131,2.000,0) -- (1.131,2.000,1.575);
  \draw[nsGuide] (0.881,0.917,0) -- (0.881,0.917,1.300);
  \draw[nsGuide] (2.000,0.825,0) -- (2.000,0.825,1.449);
  \draw[nsGuide] (1.983,0.820,0) -- (1.983,0.820,1.446);
  \draw[nsGuide] (0.888,0.911,0) -- (0.888,0.911,1.300);
  \draw[nsGuide] (0.000,0.615,0) -- (0.000,0.615,1.532);
  \draw[nsGuide] (1.233,0.308,0) -- (1.233,0.308,1.375);
  \draw[nsGuide] (1.200,0.000,0) -- (1.200,0.000,1.452);
  \filldraw[nsOrange!55, draw=black, line width=0.4pt, fill opacity=0.9] (0.881,0.917,1.300) -- (1.131,2.000,1.575) -- (0.000,2.000,2.005) -- (0.000,0.615,1.532) -- cycle;
  \filldraw[nsBlue!55, draw=black, line width=0.4pt, fill opacity=0.9] (2.000,0.825,1.449) -- (2.000,2.000,1.706) -- (1.131,2.000,1.575) -- (0.881,0.917,1.300) -- (0.888,0.911,1.300) -- (1.983,0.820,1.446) -- cycle;
  \filldraw[nsPurple!55, draw=black, line width=0.4pt, fill opacity=0.9] (1.233,0.308,1.375) -- (1.983,0.820,1.446) -- (0.888,0.911,1.300) -- cycle;
  \filldraw[nsRed!55, draw=black, line width=0.4pt, fill opacity=0.9] (0.000,0.000,1.673) -- (1.200,0.000,1.452) -- (1.233,0.308,1.375) -- (0.888,0.911,1.300) -- (0.881,0.917,1.300) -- (0.000,0.615,1.532) -- cycle;
  \filldraw[nsGreen!55, draw=black, line width=0.4pt, fill opacity=0.9] (1.200,0.000,1.452) -- (2.000,0.000,1.680) -- (2.000,0.825,1.449) -- (1.983,0.820,1.446) -- (1.233,0.308,1.375) -- cycle;
  \draw[nsRidge] (1.131,2.000,1.575) -- (0.881,0.917,1.300);
  \draw[nsRidge] (2.000,0.825,1.449) -- (1.983,0.820,1.446);
  \draw[nsRidge] (0.888,0.911,1.300) -- (1.983,0.820,1.446);
  \draw[nsRidge] (0.881,0.917,1.300) -- (0.888,0.911,1.300);
  \draw[nsRidge] (0.881,0.917,1.300) -- (0.000,0.615,1.532);
  \draw[nsRidge] (1.983,0.820,1.446) -- (1.233,0.308,1.375);
  \draw[nsRidge] (1.200,0.000,1.452) -- (1.233,0.308,1.375);
  \draw[nsRidge] (1.233,0.308,1.375) -- (0.888,0.911,1.300);
\end{tikzpicture}
}
\end{wrapfigure}
since different weighted sites can generate the same power diagram.
This redundancy is best seen by reparametrizing each weighted site $(\qq_i, \weight_i)$ as a lifted point $\smash{(\qq_i, \height_i) \in \R^{d+1}}$ with {height} $\smash{\height_i \!=\! \norm{\qq_i}^2 \!-\! \weight_i}$ \cite{Aurenhammer1987PowerDiagrams}.
In terms of lifted points, the power distance can be rewritten as $\smash{\norm{\xx}^2 \!-\! ( 2\,\qq_i\tr\xx \!-\! \height_i )}$, so each cell $\cell_i$ in \eqref{eq:power-cell} collects the points from $\domain$ where the affine function $\smash{2\,\qq_i\tr\xx \!-\! \height_i}$ is maximal.
% 
Importantly, these cells remain unchanged when all affine functions are scaled by a common factor $s \!>\! 0$ and offset by a common affine function $\smash{2\,\translation\tr\xx - c}$.
% We refer to the positions of the weighted sites that generate a power diagram as its \emph{generators}; they coincide with the sites $\qq$ here, but will differ from them once we introduce shifted sites in \S\ref{sec:contributions}.
Therefore, the weighted sites that generate the same power diagram as $(\qq, \weight)$ correspond to homotheties of the lifted points \cite{Aurenhammer1987Recognising}, denoted by the group
\begin{equation}
  \homotheties(\qq, \weight) = 
  \Bigl\{ 
    \bigl\{ s\,(\qq_i, \height_i) + (\translation, c) \bigr\}_{i=1}^n 
    \;\text{with}\; 
    s \!\in\! \R^{+},\ \translation \!\in\! \R^d,\ c \!\in\! \R 
  \Bigr\}.
  \label{eq:lifted-homothety}
\end{equation}
% where the same $(s, \translation, c)$ applies to all lifted points , Memari2012PrimalDual}.
In terms of weighted sites, each element of $\homotheties(\qq, \weight)$ has sites 
$\smash{s\,\qq_i + \translation}$ 
and weights 
$\smash{s\,\weight_i + (s^2 \!-\! s)\norm{\qq_i}^2 + 2s\,\translation\tr\qq_i + \norm{\translation}^2 - c}$, 
whose quadra\-tic dependence on the sites makes the lifted form \eqref{eq:lifted-homothety} more convenient.
This redundancy will be crucial to our analysis in \S\ref{sec:contributions}, where the $(d\!+\!2)$ infinitesimal homotheties generating $\homotheties$ always lie in, and generically span, the null-space of the Hessian of our joint formulation.

\paragraph{Energy.}
Based on \eqref{eq:power-cell}, the power diagram for given weighted sites $(\qq,\weight)$ can be interpreted as the partition of $\domain$ that minimizes the total power distance among all its partitions, with optimal value
\begin{equation}
  \energy(\qq, \weight) = \min_{\cell} \ \sum\nolimits_{i} \int_{\cell_i} \bigl( \pdist{\xx}{i} \bigr) \, d\xx .
  \label{eq:power-diagram}
\end{equation}
% We refer to $\energy$ as the \emph{energy} of the power diagram.
For zero weights, $\energy$ is often referred to as the \emph{quantization error} between the cells and their respective sites \cite{Lloyd1982Quantization}. %, and we use this term more generally for the objective of \eqref{eq:power-diagram} with zero weights at any partition.
Note that, for any fixed partition, the objective in \eqref{eq:power-diagram} is affine in the weights, so $\energy$ is a minimum of affine functions and is thus concave in $\weight$.
Moreover, 
% by Reynolds' transport theorem, the boundary terms cancel across shared facets, where neighboring power distances are equal, so 
the first derivatives of $\energy$ reduce to those of the integrand (see \eg \cite[\S2.4--2.5]{deGoes2014Thesis}):
\begin{subequations}
  \label{eq:energy-gradient}
  \begin{empheq}[left={\empheqlbrace\,}]{align}
    \partial_{\weight_i} \energy &= -\cellvol_i , \label{eq:weight-gradient} 
    \\[2mm]
    \partial_{\qq_i} \energy &= 2 \int_{\cell_i} (\qq_i - \xx) \, d\xx = 2\,\cellvol_i\,(\qq_i - \centroid_i) . \label{eq:centroid-gradient}
  \end{empheq}
\end{subequations}
In particular, stationary sites coincide with the centroids of their cells, \ie $\smash{\qq_i = \centroid_i}$, which is often referred to as a \emph{centroidal} power diagram \cite{Liu2009CVT, deGoes2012BlueNoise, Xin2016CentroidalPower}.
% As we will see in \S\ref{sec:contributions}, these gradients will serve as the building blocks of our new derivations.

\paragraph{Volume constraints.}
As detailed by \citet{Aurenhammer1998Minkowski}, power diagrams are central to the solution of the semi-discrete optimal transportation problem, which seeks the partition $\cell$ that minimizes the transportation cost from the uniform measure on $\domain$ to a discrete measure on each site $\qq_i$ with prescribed volume $\pvol_i$.
More concretely, this transportation optimization is written as
\begin{subequations}
  \label{eq:volume-constrained}
  \begin{alignat}{2}
    & \min_{\cell} \quad && \sum\nolimits_{i} \int_{\cell_i} \norm{\xx - \qq_i}^2 \, d\xx \label{eq:volume-constrained-obj} \\
    & \text{s.t.} \quad && \cellvol_i = \pvol_i \ \ \forall i . \label{eq:volume-constrained-cons}
  \end{alignat}
\end{subequations}
% 
Introducing multipliers $\weight$ for the volume constraints turns \eqref{eq:volume-constrained} into a $\max\nolimits_{\weight} \min_{\cell}$ problem.
Notably, for fixed multipliers, the inner minimization over partitions is exactly \eqref{eq:power-diagram}, thus generating a power diagram with the multipliers playing the role of weights.
Adding back the constraint terms \eqref{eq:volume-constrained-cons} yields the \emph{dual objective}
\begin{equation}
  \dualenergy(\weight)
    = \energy(\qq, \weight) + \sum\nolimits_{i} \weight_i \, \pvol_i.
  \label{eq:volume-dual}
\end{equation}
% where only the weights $\weight$ are optimization variables, while the sites $\qq$ and volumes $\pvol$ are fixed data.
% For all dual objectives throughout our presentation, we list only the optimization variables as arguments and leave the fixed data implicit from the context.
% 
The volume-constrained problem in \eqref{eq:volume-constrained} thus reduces to the {unconstrained} maximization of $\dualenergy(\weight)$ over the space of power diagrams, parametri\-zed by the weights $\weight$.
By \eqref{eq:weight-gradient}, the gradient of $\dualenergy$ recovers the volume constraint
\begin{equation}
  \partial_{\weight_i} \dualenergy = \pvol_i - \cellvol_i ,
  \label{eq:volume-dual-gradient}
\end{equation}
so stationary weights are those whose cells realize the prescribed volumes.
As derived by \citet{deGoes2012BlueNoise}, the Hessian of $\dualenergy$ is a (negated) Laplacian matrix on the power diagram connectivity with entries
\begin{subequations}
  \label{eq:newton-step}
  \begin{empheq}[left={\empheqlbrace\,}]{align}
    \partial^2_{\weight_i \weight_j} \dualenergy
      &= \frac{1}{2} \frac{\area_{ij}}{\norm{\qq_i - \qq_j}},
    \label{eq:newton-step-off}
    \\[2mm]
    \partial^2_{\weight_i \weight_i} \dualenergy
      &= -\!\sum\nolimits_{j} \partial^2_{\weight_i \weight_j} \dualenergy .
    \label{eq:newton-step-diag}
  \end{empheq}
\end{subequations}
Therefore, this Hessian is negative semi-definite with the constant vector in its kernel, reflecting that adding a common constant to all weights leaves the power diagram unchanged.

\paragraph{Achievable centroids.}
Finally, we summarize results by \citet{Bourne2025Inverting, Bourne2026ConvexOrder} that describe the set of centroids achievable by all partitions of $\domain$ with prescribed volumes.
This set, which we denote as $\achievable$, is convex and compact, contained within the affine subspace defined by identity \eqref{eq:partition-feasibility-centroid} \cite[\S3.1]{Bourne2026ConvexOrder}.
% 
Volume-constrained power diagrams play a special role in $\achievable$: the centroids of power diagrams with the prescribed volumes are the exposed points of $\achievable$, \ie the points where a supporting hyperplane touches $\achievable$ and nowhere else \cite[Thm.~3.6]{Bourne2026ConvexOrder}.
Moreover, the sites of a power diagram are an outward normal of $\achievable$ at its centroids \cite[Prop.~3.5]{Bourne2026ConvexOrder}, and so are the sites of every element of $\homotheties$ in \eqref{eq:lifted-homothety}.
% , since scaling preserves the direction of the normal, while translation only adds a component orthogonal to the subspace of \eqref{eq:partition-feasibility-centroid} (in the volume-weighted $\smash{L^2}$ sense).
These homotheties form the {normal cone} of $\achievable$ at these centroids, \ie the set of all outward normals of supporting hyperplanes, since each of them generates the unique power diagram with the same volumes and centroids \cite[Thm.~4.5]{Bourne2025Inverting}.
Therefore, power diagrams lie on the boundary of $\achievable$ relative to \eqref{eq:partition-feasibility-centroid}.
The rest of this boundary consists of flat faces exposed by coinciding sites, while its interior includes only centroids of partitions that are not power diagrams.
Based on these results, we call centroids \emph{achievable} when they belong to $\achievable$, and \emph{realizable} when they are the centroids of a power diagram, while we call centroids \emph{unachievable} when they verify identity \eqref{eq:partition-feasibility-centroid} but do not belong to $\achievable$, \ie when no partition of $\domain$ with volumes $\pvol$ has them as centroids. % \ie the exposed points of $\achievable$.
%
We rely on this characterization to determine which power diagrams our formulations return, as detailed next.
% 
% As we construct power diagrams with prescribed centroids next in \S\ref{sec:contributions}, this geometry will help us determine when such diagrams exist, and which diagram our formulations return otherwise.

\end{document}
